% Tier 1 parity doc: sectioning, ToC, cross-refs, footnotes, hyperref, bib.
% Exercises multi-pass label resolution, marks, and whatsit emission.
\documentclass[11pt]{article}
\usepackage{amsmath}
\usepackage{hyperref}
\begin{document}
\tableofcontents

\section{Introduction}\label{sec:intro}
This document references Section~\ref{sec:method} on
page~\pageref{sec:method}, and cites~\cite{knuth1984}.
Here is a footnote.\footnote{Footnotes exercise insertion nodes.}

\section{Method}\label{sec:method}
As discussed in Section~\ref{sec:intro}, we proceed. A hyperlink to
\href{https://example.org}{an example site} and an internal jump to
Section~\ref{sec:intro}.

\section{Conclusion}
See equation~\eqref{eq:e} below.
\begin{equation}\label{eq:e}
  e = \lim_{n\to\infty}\left(1+\tfrac{1}{n}\right)^n
\end{equation}

\begin{thebibliography}{9}
\bibitem{knuth1984} D.~E. Knuth, \emph{The \TeX book}, Addison-Wesley, 1984.
\end{thebibliography}

\end{document}
